{\bfseries Transformada de Fourier cuántica (QFT). Considere un registro de $n = 2$ qubits (estados base $\ket{00}$, $\ket{01}$, $\ket{10}$, $\ket{11}$). Aplique la QFT al estado base $\ket{x}$ con $x = 2$ (binario $10$). Use la representación producto:}

\[
	\text{QFT}\ket{x_1 x_2} = \frac{1}{2}\left(\ket{0} + e^{2\pi i(0.x_1 x_2)}\ket{1}\right) \otimes \left(\ket{0} + e^{2\pi i(0.x_2)}\ket{1}\right),
\]

{\bfseries donde $x_1$ es el bit más significativo.}

\begin{enumerate}[a)]
	\item {\bfseries Escriba los bits $x_1$, $x_2$ de $x = 2$.}

	\vspace{3cm}

	\item {\bfseries Calcule las fracciones binarias $0.x_1 x_2$ y $0.x_2$ en forma decimal (o fraccionaria).}

	\vspace{3cm}

	\item {\bfseries Determine los factores de fase $e^{2\pi i(0.x_1 x_2)}$ y $e^{2\pi i(0.x_2)}$.}

	\vspace{3cm}

	\item {\bfseries Escriba explícitamente el estado $\text{QFT}\ket{10}$ como producto de dos qubits y también como combinación lineal de los cuatro estados base $\ket{00}$, $\ket{01}$, $\ket{10}$, $\ket{11}$. Verifique que el resultado coincide con el obtenido mediante la definición matricial de la QFT para $n = 2$.}

	\vspace{5cm}

\end{enumerate}
